\label{night:P3F.46}
\StatementRecord{night:P3F.46}{proved}{}{\cref{night:P3F.35}, \cref{iv:correction:P3K.22}}
\par\noindent\textit{Dependency scope.} We use \cref{iv:correction:P3K.22} only at their stated retained-construction or corrected-conclusion scopes.


Let $S$ be the minimal regular surface of
\[
 y^2=x^3+tx+t^2,\qquad G=(0,t),\qquad
 M=\mathcal O_S(aG-O),\qquad Q=bG.
\]
Use its $I_1$ fibre at $t_c=-4/27$ as the sole simple whole-fibre
zero. The following four pairs admit the actual two-cap families
constructed below:
\[
\begin{array}{c|c|c}
(a,b)&m_6\text{ at }IV^*&m_2\text{ at }III\\ \hline
(1,6)&3l,\ l\ge1\text{ (good)}&4k,\ k\ge1\text{ (good)}\\
(2,3)&3l,\ l\ge1\text{ (good)}&4k+1,\ k\ge0\text{ (reduced)}\\
(3,2)&3l+2,\ l\ge0\text{ (reduced)}&4k,\ k\ge1\text{ (good)}\\
(6,1)&3l+2,\ l\ge0\text{ (reduced)}&4k+1,\ k\ge0\text{ (reduced)}.
\end{array}                                                     \tag{F46.1}
\]
For each pair choose signed integers $n_i$ with
$\gcd(m_i,n_i)=1$. The marked meridians are
$m_i\alpha_i+n_i\gamma$, where $\gamma$ is the outward
suspension generator. With $L=\operatorname{lcm}(m_6,m_2)$,
the construction is a proper holomorphic family over a punctured
$r$-disc with scalar $\lambda=r^L$. Its fibres are the standard
oriented $S^6$ exactly when
\[
                  D=m_2n_6+m_6n_2\in\{1,-1\}.                 \tag{F46.2}
\]
This criterion holds for every allowed fan. The statement assigns
the integral topology and retains the actual parameter closing;
it assigns no parameter-return class.
\end{theorem}
\begin{proof}
\begin{enumerate}[label=\arabic*.]
\item \par\medskip\noindent\textbf{The source and the exact root maps.}
The discriminant is $t^3(4+27t)$, so the fibres are $III$ at zero,
$I_1$ at $t_c$, and $IV^*$ at infinity. On $x\ne0$ put
$u=(y-t)/x$ and $v=(y+t)/x$. The inverse formulas are
\[
 x=uv-(v-u)/2,\quad t=x(v-u)/2,\quad y=x(u+v)/2.
\]
Thus $S$ is rational. The marked section meets nonidentity
components in $E_6$ and $A_1$, has $G\cdot O=0$, and has
height $1/6$. The integral source Gram matrix below has determinant
$-1$; the discriminant and height identity give primitive $G$
and no Mordell--Weil torsion.

\item The source map is $h=t_{-bG}$. Write $T$ for its pushforward on
divisor cycles. The exact section formula below gives
$TM-M=-F$. Its pullback is $T^{-1}$, so $h^*M-M=F$.
Consequently the forward line map $M_p\to M_{h(p)}$ is a section
of $h^*M\otimes M^{-1}=\mathcal O(F)$. Choose its section with
the sole simple zero at $t_c$, denote it by $B_0$, and put
$B_\lambda=\lambda B_0$. This specifies the line direction.

\item At infinity set $q=1/t$, $x_0=q^2x$, $y_0=q^3y$. Then
\[
 y_0^2=x_0^3+q^3x_0+q^4,\qquad G=(0,q^2).
\]
For $m_6=3l+2$, substitute
$q=z^{m_6}$, $x_0=z^{4l+2}X$, $y_0=z^{6l+3}Y$. We obtain
\[
 Y^2=X^3+z^{l+2}X+z^2,\qquad G'=(0,z).                         \tag{F46.3}
\]
Its $A_2$ resolution has a reduced $IV$ fibre. In
$X=zX_1,Y=zY_1$, the nonidentity labels are $Y_1=1,-1$.
For $k=2l+1$ the actual deck is
\[
 D(X,Y,z)=(\zeta_m^{-2k}X,\zeta_m^{-3k}Y,\zeta_m z).
\]
Its character on $Y_1$ is
$\zeta_m^{-(3k+1)}=\zeta_m^{-2m}=1$, so it fixes all
three labels. For $m_6=3l$, the substitution $k=2l$ instead gives
\[
 Y^2=X^3+z^lX+1,\qquad G'=(0,1),                               \tag{F46.4}
\]
with central deck $(X,Y)\mapsto(\zeta_3^2X,Y)$.
The point $(0,1)$ has exact order three: its horizontal tangent
has triple intersection with $Y^2=X^3+1$.

\item At the $III$ site take $q=t$. For $m_2=4k+1$, put
$q=z^{m_2}$, $x=z^{2k}X$, $y=z^{3k}Y$. The equation is
\[
 Y^2=X^3+zX+z^{2k+2},\qquad G'=(0,z^{k+1}).                    \tag{F46.5}
\]
The $A_1$ resolution gives the reduced $III$ fibre, whose two
labels are fixed by the deck. For $m_2=4k$ the same substitution
gives
\[
 Y^2=X^3+X+z^{2k},\qquad G'=(0,z^k).                           \tag{F46.6}
\]
Its central deck is $(X,Y)\mapsto(-X,iY)$ and its marked
limit $(0,0)$ is the nonzero fixed two-torsion point.
These substitutions and the stated resolutions are the actual
maps to the source models.

\par\medskip\noindent\textbf{The two line comparisons.}
At a reduced root, $aG'$ has identity component and
$N=\mathcal O(aG'-O)$ has degree zero on every component.
The original $M$ has the same zero multidegrees at that site.
Pull it to the resolved fibre product first. Every curve
contracted in relative minimalization has degree zero, so the
blowup Picard decomposition descends the line through each
actual contraction. Its comparison with $N$ is a vertical
divisor with zero component degrees. The reduced $III$ or $IV$
intersection matrix has integral kernel the whole-fibre vector.
The corresponding base monomial therefore gives the line
isomorphism and retains its deck character and the specified $B_0$.

\item For a good $IV^*$ root, the minimal cubic deck has tangent
weights $(2,1)$ at its three fixed points. Its quotient lattice
is $\mathbb Z^2+\mathbb Z(2,1)/3$, and the minimal quotient
rays are $((3-j)/3,j/3)$. A common smooth root refinement has rays
\[
                  (1,0),(2,1),(1,1),(1,2),(0,1).
\]
Let $\beta$ be this root blowup and $\pi$ its toric map to the
minimal quotient resolution. The section $w=0$ has Cartier
support $h(v)=\max(v_w-2v_z,0)$ there. Pulling back this
function gives the exact section-divisor difference
\[
 \beta^*(\text{root section})-\pi^*(\text{resolved section})
       =\sum_v\min(v_w,2v_z)E_v,
\]
with coefficients $(2,1,1)$. Subtract the $O$ contribution from
the $aG'$ contribution to obtain
\[
                    \beta^*N=\pi^*M+D_{aG}-D_O.                \tag{F46.7}
\]
The invariant Weierstrass substitution is the source map, and
its crepant quotient resolution is the original $IV^*$ model.
Since $Q'_0=O$ for $b=3,6$, its translation is locally
$w\mapsto w\,u+z^2h(z^3)$. It preserves
$(w,z),(w,z^2),(z,w),(z,w^2)$, hence the common normalized
blowup and every correction in (F46.7). For $m_6=3l$, pull this
entire correspondence through the extra degree-$l$ root and
resolve the fibre product. The divisor equality and maps pull
back exactly.

\item For a good $III$ root, the minimal fourth deck has two
order-four fixed points $O,P$ and one orbit $\{R_1,R_2\}$
with stabilizer two. The tangent weights are $(1,1)$.
Blow up these four points by $\beta$ and take the quotient
$\pi$. Its exceptional images $D_0,D_1,D_2$ have squares
$-4,-4,-2$, and the central image $C$ has square $-1$.
The actual fibre is $4C+D_0+D_1+2D_2$. Contract $C$, then the
resulting $(-1)$-curve $D_2$. The remaining curves have squares
$-2$ and intersection two; the invariant Weierstrass map gives
their original $III$ tangency. The marked sections avoid both
contracted curves. Thus, for these contractions $\alpha$,
\[
                 \beta^*N=\pi^*\alpha^*M+E_{aG}-E_O.           \tag{F46.8}
\]
For $b=2,6$ we have $Q'_0=O$, so its translation fixes all four
centres and preserves the correction. This carries the actual
$B_0$ and its deck commutation to $N$. Higher fourth roots
pull back this whole diagram and its characters.

\item The good central line has nonzero torsion Picard class of
order three at $IV^*$ and order two at $III$. The actual map
\[
 (\mathbb C/\Lambda)^D\longrightarrow\Lambda/(D-1)\Lambda,
                \quad[v]\longmapsto[(D-1)v]
\]
carries it to the primitive torsion Euler class of the cap
circle line. A common deck character is retained and does not
change that torsion order.

\par\medskip\noindent\textbf{The exact norms and the joint family.}
At either site write $m,n,N,D$ for its data. Multiplication by
$m$ is invertible on the reduced additive identity group and
component group. At a good root use its analytic inverse branch
at $O$. Thus the unique marked analytic division $mR=Q'$ is
$D$-invariant. The source map $t_{-nR}D$ has $m$th power
$t_{-nQ'}$. It extends on the minimal regular root: resolve the
addition graph and compare the invariant holomorphic two-form;
the common exceptional support gives the holomorphic map and
its inverse.

\item Lift $t_{-R}$ to $F$ on $N$. The fibre chord divisor gives its
Hom section, and its residual vertical divisor has zero
component degrees. Remove the resulting whole-fibre base power.
Finite deck averaging then gives
$DFD^{-1}=\chi F$, with $\chi^m=1$.
The specified descended lift is $B_0=\beta(q)F^m$.
The ratio is a base unit because the fibres are compact and
connected, and deck commutation makes it invariant. Hence
\[
       (F^nD)^m/B_0^n
          =\chi^{nm(m-1)/2}\beta(q)^{-n}.
\]
Divide by an invariant analytic $m$th root of this unit.
The resulting $A_0$ satisfies $A_0^m=B_0^n$ and
$[A_0,B_0]=1$, retaining the character factor for even $m$.

\item Set $\lambda=r^L$ and
$A_{i,r}=r^{Ln_i/m_i}A_{i,0}$. For integers $u_i,v_i$ with
$n_i u_i+m_i v_i=1$, put
\[
 H_i=A_{i,r}^{u_i}B_\lambda^{v_i}
       =r^{L/m_i}H_{i,0},\qquad
 H_i^{m_i}=B_\lambda,\quad H_i^{n_i}=A_{i,r}.                  \tag{F46.9}
\]
On smaller compact source tubes the norm ratios of $H_{i,0}$
and its inverse are bounded. Thus $H_i$ is strictly contracting
for small $r$, and its norm shell gives a smooth proper cap
$N_i^\times/\langle H_i\rangle$. A unit rotation equal to
$H_i^j$ forces $j=0$ by contraction, then forces the unit to
be one. The delta circle is therefore free on both caps.
For outward $h_i=H_i^{-1}$, the single-$A$ cut gives exactly
\[
 \alpha_i\longmapsto-n_i h_i,\qquad
 \gamma\longmapsto m_i h_i,\qquad
 \mu_i=m_i\alpha_i+n_i\gamma.                                 \tag{F46.10}
\]

\item The holomorphic annular map is constructed in the actual
constant-period extension $(\mathbb C\times\mathbb C^*)/\Lambda$.
Remove its integral period two-cocycle on the annulus, whose
$H^2$ is zero, then solve the holomorphic Lie-bundle cocycle
on the Stein annulus. Choose the gamma component whose positive
cut is $A_{i,r}$ itself. For $c=\log\lambda/(2\pi i)$ its scalar
variation is $n_i(c-c_0)s_i\delta/m_i$. This is (F46.10) with
every elliptic and delta phase retained.

\item For a positive $r$-loop put $k_i=Ln_i/m_i$. The inverse closing
near cap $i$ is $q_i^{k_i}$ times its extending unit. On the
ordinary region the exact global multiplier is
\[
                  f(t)=t^{k_2}(t-t_c)^{-(k_6+k_2)}.              \tag{F46.11}
\]
Its orders are $k_2$ at zero, $k_6$ at infinity, and
$-(k_6+k_2)$ at the selected cusp. Its ratios to the two local
multipliers are units. The cap charts close by the identity,
and their annular maps supply these multipliers. At the cusp,
(F46.11) is, modulo the specified $B_\lambda$ iterates, a bounded
source translation and an invertible scalar unit. It and its
inverse extend in the actual nodal Mumford charts. This proves
holomorphic parameter closing. Compact norm shells give
properness on each compact parameter annulus.

\par\medskip\noindent\textbf{The full source and boundary complexes.}
Use the integral basis $O,F,C_0,\ldots,C_5,A_1,G$.
The $E_6$ graph has edges
$(0,1),(1,2),(2,3),(3,4),(2,5)$, Cartan matrix $A$, and
highest-root marks $v=(1,2,3,2,1,2)^t$.
The source Gram matrix $B$ has $O^2=G^2=-1$, $OF=GF=1$,
root blocks $-A$ and $(-2)$, and $GC_0=GA_1=1$;
all unspecified pairings vanish. Its determinant is $-1$,
so these are the entire integral source lattice.

\item For an integer $j$, let $R_j$ be the inverse-Cartan column at
$C_0$ or $C_4$ when $j=1$ or $2$ modulo three, together
with $A_1/2$ when $j$ is odd. Let $c_j$ be the sum of the
corresponding diagonal entries $4/3$ and $1/2$.
The actual section formula is
\[
 \phi=G-O-F+R_1,\quad \phi^2=-1/6,\qquad
 P_j=O+\frac{j^2/6+c_j}{2}F+j\phi-R_j.                          \tag{F46.12}
\]
It is integral, has square $-1$, and has the prescribed component
intersections. Translation by $G$ sends $O$ to $G$, $G$ to
$P_2$, cyclically permutes the three affine $E_6$ arms, and
swaps $A_1$ with $F-A_1$. These actual columns define the
integral isometry $T_+$, and $T=T_+^{-b}$.

\item Put
\[
 L_6=(F-\textstyle\sum v_jC_j,C_0,\ldots,C_5),\quad
 L_2=(F-A_1,A_1),\quad L=(L_6,L_2).
\]
Their multiplicity vectors are $(1,v)^t$ and $(1,1)^t$;
both products with $L_i$ equal $F$. If $P$ is the actual
component permutation, then $TL=LP$.
For the complement $V$ of the two fibre tubes the integral
source/relative complex is
\[
\begin{gathered}
 C_4=\mathbb Z,\quad C_3=\mathbb Z^2,\quad C_2=\operatorname{NS}(S),\\
 C_1=\mathbb Z^9,\quad C_0=\mathbb Z,\\
 d_4=(1,1)^t,\quad d_2=L^tB,\quad d_3=d_1=0.
\end{gathered}                  \tag{F46.13}
\]
Boundary $i$ has $C_3=\mathbb Z$, $C_2=C_1=\mathbb Z^{r_i}$,
$C_0=\mathbb Z$, and $d_2=L_i^tBL_i$.
Its inclusion is $L_i$ in degree two, the component-coordinate
embedding in degree one, and the $i$th coordinate in degree
three. These are the complexes of the actual pair. The source
and the tube retracts have free even homology, so the degree-one
derived homotopy group of their even split complexes is zero.
The displayed maps therefore have no omitted odd homotopy
parameter. This argument includes the actual tangencies.

\item The source action is $T$ in degree two, $P^{-t}$ in degree one,
and identity otherwise. The boundary action is $P_i$ in degree
two, $P_i^{-t}$ in degree one, and identity otherwise.
The program forms both actual mapping-torus cones, with
\[
       d_{{\rm cone},j}=
       \begin{pmatrix}d_j&T_{j-1}-I\\0&-d_{j-1}\end{pmatrix},
\]
and their block-diagonal inclusions. All chain matrices,
integral Smith representatives, projections, and marked
$\alpha,\gamma,N_3,F\gamma$ cycles are exported. They give
\[
 H_1(W)=\mathbb Z^2,\quad H_3(W)=\mathbb Z^2,\quad
 H_4(W)=\mathbb Z,\quad
 H_2(W)=\mathbb Z/b\oplus\mathbb Z^2,                           \tag{F46.14}
\]
with the trivial summand omitted when $b=1$.

\item We also compute the full fundamental group. The remaining
$I_1$ kills its primitive vanishing cycle $v$. If $A$ is the
$III$ monodromy, then $A^2=-I$ and $\operatorname{tr}A=0$.
The product relation with $IV^*$ and this $I_1$ gives
$|\langle v,Av\rangle|=1$: expand
$\operatorname{tr}(A^{-1}T_v^{-1})$, whose absolute value is
that pairing and is one from the inverse $IV^*$ trace.
Thus $v,Av$ form an integral elliptic basis. Conjugating the
vanishing relation by the retained base loop kills both.
The section gives the remaining base-annulus generator, so
$\pi_1(V)=\mathbb Z$. Translation fixes it and
\[
             \pi_1(W)=\mathbb Z\alpha\oplus\mathbb Z\gamma.
\]
The two boundary base meridians map to $\alpha,-\alpha$;
both outward suspension loops map to $\gamma$.

\par\medskip\noindent\textbf{The cap specializations.}
A good cap has
$H_1=\mathbb Zh\oplus\mathbb Z/p$, $H_2=\mathbb ZC$,
$H_3=\mathbb Z(Ch)$, $H_4=0$, with $p=3$ or two.
Actual smooth-root transport sends the generic fibre to its
central elliptic fibre with degree one. The same transport is
the identity on the elliptic lattice. If $nu+mv=1$, then $u$
is prime to the deck order and
\[
                    (D^u-I)\Lambda=(D-I)\Lambda.
\]
For order three this follows from $D+I=-D^2$; for order four
the relevant odd powers are $D$ or $D^{-1}$.
The multiplying matrices are unimodular. Thus the incoming
finite torsion maps isomorphically to cap torsion, and the
incoming full elliptic group maps onto the cap elliptic lattice.

\item A reduced cap has $H_1=\mathbb Zh$, $H_2=\mathbb ZC$,
$H_3=\mathbb Z(\sum C_j)$, $H_4=0$.
Its primitive $H_i$ cyclically permutes all central components:
its component translation is $-Q/m_i$, independent of $n_i$
by (F46.9), and is nonzero modulo three or two.
The actual fibre specializes to $\sum C_j$. Its image in
coinvariants is $r_iC$, while its invariant trace is primitive.

\item Use boundary $H_2$ coordinates $(\eta_i,r_i,s_i)$ in the good
case, with $\eta_i$ of order $p_i$ and $F=r_i$.
In the reduced case use $(r_i,s_i)$, with $F=3r_i$ at $IV^*$
and $F=2r_i$ at $III$. The actual cap maps $r_i\mapsto C_i$.
Write its remaining column as $s_i\mapsto\kappa_iC_i$ and
retain that integer.
The full exterior inclusions, with torsion row first, are
\[
\begin{array}{c|c|c}
(a,b)&IV^*&III\\ \hline
(1,6)&\begin{pmatrix}-2&1&6\\0&0&-1\\0&0&-1\end{pmatrix}
 &\begin{pmatrix}3&1&3\\0&0&0\\0&0&1\end{pmatrix}\\[5pt]
(2,3)&\begin{pmatrix}1&2&7\\0&0&0\\0&0&-1\end{pmatrix}
 &\begin{pmatrix}1&0\\0&-1\\0&2\end{pmatrix}\\[5pt]
(3,2)&\begin{pmatrix}-1&18\\0&1\\0&3\end{pmatrix}
 &\begin{pmatrix}-1&-3&-1\\0&0&0\\0&0&1\end{pmatrix}\\[5pt]
(6,1)&\begin{pmatrix}0&-1\\0&3\end{pmatrix}
 &\begin{pmatrix}0&-1\\0&2\end{pmatrix}.
\end{array}                                                     \tag{F46.15}
\]
The first row is taken modulo $b$ when $b>1$.
For $b=6$ the exact CRT map is
$\mathbb Z/6\to\mathbb Z/3\oplus\mathbb Z/2$,
$1\mapsto(\eta_6,\eta_2)$, with exterior coefficient
$-2+3=1$. For $b=3$ the torsion generator already maps to one.
For $b=2$ use minus the incoming $III$ torsion generator.

\item In domain order $(\eta,r_6,s_6,r_2,s_2)$ and target order
$(\tau,g_1,g_2,C_6,C_2)$, the complete degree-two matrices are
\[
 D_{2,6}=\begin{pmatrix}
1&1&0&1&3\\0&0&-1&0&0\\0&0&-1&0&1\\
0&-1&-\kappa_6&0&0\\0&0&0&-1&-\kappa_2
\end{pmatrix},\quad
 D_{2,3}=\begin{pmatrix}
1&2&1&1&0\\0&0&0&0&-1\\0&0&-1&0&2\\
0&-1&-\kappa_6&0&0\\0&0&0&-1&-\kappa_2
\end{pmatrix}.                                                  \tag{F46.16}
\]
\[
 D_{2,2}=\begin{pmatrix}
1&1&0&1&1\\0&0&1&0&0\\0&0&3&0&1\\
0&-1&-\kappa_6&0&0\\0&0&0&-1&-\kappa_2
\end{pmatrix},\qquad
 D_{2,1}=\begin{pmatrix}
0&-1&0&-1\\0&3&0&2\\-1&-\kappa_6&0&0\\
0&0&-1&-\kappa_2
\end{pmatrix}.                                                  \tag{F46.17}
\]
Their determinants are $1,1,-1,-1$. For $b>1$, the first
coordinate and row have order $b$; the first column has no
free entry. Thus these are integral group isomorphisms for
every value of both actual collar coefficients.

\par\medskip\noindent\textbf{All remaining gluing degrees.}
In degree one every good boundary torsion maps isomorphically
to cap torsion and to zero in $W$. The exported full matrices
retain arbitrary torsion path coefficients $c_i,d_i$ in the
images of $\alpha_i,\gamma$. Eliminating precisely those finite
isomorphisms gives, in physical order
$(\alpha_6,\gamma_6,\alpha_2,\gamma_2)$,
\[
 D_1=\begin{pmatrix}
1&0&-1&0\\0&1&0&1\\n_6&-m_6&0&0\\0&0&n_2&-m_2
\end{pmatrix}.                                                  \tag{F46.18}
\]
In degree three use the actual oriented source boundary $N_i$
followed by $F\gamma_i$. The cap map is $(-n_i,m_i)$:
the norm-circle degrees in (F46.10) are followed by specialization
to the primitive invariant central fibre. The cap source
retracts to its two-dimensional central fibre, so it contributes
no additional three-class. The exterior has $N_6=-N_2$ and
the same $F\gamma$ at both ends. In the exported Smith bases
the $IV^*$ boundary change is $\operatorname{diag}(1,-1)$,
the $III$ change is identity, and the exterior target change
for $b=1,2$ is $\left(\begin{smallmatrix}0&-1\\1&0\end{smallmatrix}\right)$;
for $b=3,6$ it is identity. Consequently the exact physical
$D_3$ equals (F46.18), and
\[
                   \det D_1=\det D_3=-D.                        \tag{F46.19}
\]
In degree four the map is $(1,-1):\mathbb Z^2\to\mathbb Z$,
up to the recorded common target orientation. In degree zero
it is the actual incidence matrix
$\left(\begin{smallmatrix}1&1\\-1&0\\0&-1\end{smallmatrix}\right)$.
This specifies every nonzero Mayer--Vietoris degree.

\par\medskip\noindent\textbf{The exact fundamental group and sphere criterion.}
The full incoming elliptic lattice dies in $V$. In each good
cap it is the full normal elliptic subgroup, so van Kampen
kills it. A reduced cap has no remaining elliptic subgroup.
The exact residual presentation is
\[
 \left\langle h_6,h_2\ \middle|\
 h_6^{m_6}=h_2^{m_2},\quad h_6^{n_6}=h_2^{-n_2}\right\rangle.
\]
Choose $u,v$ with $n_6u+m_6v=1$. Then
\[
 h_6=(h_6^{n_6})^u(h_6^{m_6})^v
       =h_2^{-n_2u+m_2v},
\]
so the group is cyclic. Its abelian relation matrix has a
primitive relation and determinant $D$, giving
$\pi_1(Y)=\mathbb Z/|D|$, with $\mathbb Z$ when $D=0$.
The full integral Mayer--Vietoris sequence gives, for $D\ne0$,
\[
 H_0(Y)=H_5(Y)=\mathbb Z,\quad
 H_1(Y)=H_3(Y)=\mathbb Z/|D|,\quad H_2(Y)=H_4(Y)=0.             \tag{F46.20}
\]
For $D=0$, the first three Smith factors of (F46.18) are one
and the last is zero; all $H_j(Y)$, $0\le j\le5$, are
$\mathbb Z$. Thus $Y$ is the standard $S^5$ exactly when
$D=1$ or $-1$.

\item Both added caps have free delta action. The only fixed locus
is the chosen nodal-fibre sphere, with the primitive
$\mathcal O(-1)^{\oplus2}$ Hopf chart. The exact Gysin and
fixed-tube maps of P3K.22 give $\pi_1(X)=\pi_1(Y)$ and, for
$D=\pm1$, the integral $S^6$ homology. Since $\Theta_6=0$,
the oriented total space is the standard sphere. For all other
$D$ its displayed nontrivial fundamental group excludes a
sphere. The actual finite flop chain of P3F.35 is the identity
outside the cusp and transports this criterion to every allowed
fan. Equation (F46.2) also forces coprimality of the two orders
and of each local norm pair.

\item The complete certificate is
\path{code/part3/iv_star_iii_joint.py}, with its saved JSON and
output. It contains every source and boundary chain matrix,
Smith representative, marked cycle, finite-group degree-one
map, symbolic collar coefficient, and degree-three orientation
change. The four degree-two determinants and the common
degree-one/three determinant are symbolic identities. A further
$1056$ exact Smith cases, including $D=0$, supplement these
all-order calculations.

\end{enumerate}
\end{proof}